% answers_short.tex —— 2010 年·普通高等学校招生考试·浙江卷·文科 · 简短答案速查
% 由 tools/make_short_answers.py 自 answers.tex 的 \gksolhead/\gkanswer 提取，
% 并与 2010/2010_数学真题_浙江文_答案要点.md 逐题比对后生成，勿手改（改源 MD 后重跑）。

\begin{gknotice}
\textbf{使用说明}：本节只给最终答案便于快速核对，完整推导见解析版。答案为编写者独立推导，\textbf{非官方标准答案}。
\end{gknotice}

\gkpart{选择题}
\begin{center}\begin{tabular}{ccccccccccc}
\toprule
题号 & 1 & 2 & 3 & 4 & 5 & 6 & 7 & 8 & 9 & 10 \\
\midrule
答案 & D & B & C & A & A & B & A & B & B & D \\
\bottomrule
\end{tabular}\end{center}

\gkpart{填空题}
\begin{itemize}[leftmargin=2.6em]
  \item[\textbf{11.}] $45$，$46$。
  \item[\textbf{12.}] $\dfrac{\pi}{2}$。
  \item[\textbf{13.}] $\sqrt{10}$。
  \item[\textbf{14.}] $n(n+1)$。
  \item[\textbf{15.}] $18$。
  \item[\textbf{16.}] $20$。
  \item[\textbf{17.}] $\dfrac34$。
\end{itemize}

\gkpart{解答题}
\begin{itemize}[leftmargin=2.6em]
  \item[\textbf{18.}] （1）$\dfrac{\pi}{3}$；（2）$\sqrt3$。
  \item[\textbf{19.}] （1）$S_6=-3$，$a_1=7$；（2）$(-\infty,-2\sqrt2]\cup[2\sqrt2,+\infty)$。
  \item[\textbf{20.}] （1）证明；（2）$\dfrac12$。
  \item[\textbf{21.}] （1）$x-y-2=0$；（2）$x_4=\dfrac{2a+b}{3}$。
  \item[\textbf{22.}] （1）$y^2=8x$；（2）证明。
\end{itemize}
